\documentclass[aspectratio=169]{beamer}
\usepackage{tikz}
\usepackage{graphicx}
\usepackage{listings}
\usepackage{xcolor}
\usepackage{hyperref}
\usepackage[spanish,es-noquoting]{babel}
\usepackage[round]{natbib}

\lstset{
  basicstyle=\ttfamily\small,
  keywordstyle=\color{blue},
  stringstyle=\color{red},
  commentstyle=\color{gray},
  showstringspaces=false
}
% Minimal look
\usetheme{Warsaw}
%\setbeamertemplate{headline}
\setbeamertemplate{footline}{
  \leavevmode%
  \hbox{%
    % Left: Section name
    \begin{beamercolorbox}[wd=.33\paperwidth,ht=2.25ex,dp=1ex,left]{author in head/foot}%
      \hspace*{2ex}\insertsection
    \end{beamercolorbox}%
    % Center: Presentation title
    \begin{beamercolorbox}[wd=.34\paperwidth,ht=2.25ex,dp=1ex,center]{title in head/foot}%
      \insertshorttitle
    \end{beamercolorbox}%
    % Right: Slide numbers
    \begin{beamercolorbox}[wd=.33\paperwidth,ht=2.25ex,dp=1ex,right]{date in head/foot}%
      \insertframenumber{} / \inserttotalframenumber\hspace*{2ex}
    \end{beamercolorbox}%
  }%
  \vskip0pt%
}

\usecolortheme{default}
\usefonttheme{serif}
\graphicspath{{../}{./}}
\setbeamertemplate{caption}[numbered]

% Title info
\title[Módulo 2: Cómo Hablar con la I]{Módulo 2: Cómo Hablar con la IA}
\subtitle{Ingeniería de Prompts en la Vida Real}
\author{Ricardo Reyes-Jimenez}
\institute{IA para Todos}
\date{\today}



\begin{document}

% Title slide
\begin{frame}
  \titlepage
\end{frame}

% Reset background for normal slides
\usebackgroundtemplate{}



\begin{frame}[fragile]{Objetivo de la Sesión}


\textbf{Al finalizar este módulo, los participantes serán capaces de:}

1. Comprender la diferencia entre un motor de búsqueda y un modelo de lenguaje.
2. Estructurar peticiones precisas usando el marco de ``Los 5 Pilares''.
3. Redactar comunicaciones efectivas aplicando el marco TAM (Tarea, Audiencia, Método).

\end{frame}



\begin{frame}[fragile]{1. El Mito del Buscador}


\textbf{La IA no es Google. No busca, sino que razona probabilísticamente \cite{openai2023}.}

\begin{itemize}
  \item \textbf{Búsqueda en Google (Palabras clave):} ``Receta de pollo al horno''
  \item \textbf{Búsqueda con IA (Contexto e Intención):} ``Actúa como un chef experto. Dame una receta de pollo al horno para 4 personas, considerando que una de ellas es alérgica al ajo y otra no come lácteos. Enumera los ingredientes y los pasos de forma clara.''
\end{itemize}
\textbf{Regla General:} Si tu petición (prompt) es más corta que un Tweet, probablemente estás subutilizando la herramienta y arriesgándote a recibir respuestas genéricas.

\end{frame}



\begin{frame}[fragile]{2. Los 5 Pilares de una Petición Exitosa}


Una petición vaga produce un resultado inútil. Una petición profesional necesita 5 pilares:

1. \textbf{Rol:} Identidad de la IA (Ej. Nutricionista, Abogado, Ingeniero).
2. \textbf{Contexto:} Hechos materiales y objetivos (Ej. Preparando un viaje, negociando un aumento).
3. \textbf{Instrucción:} El verbo imperativo (Analizar, Extraer, Comparar).
4. \textbf{Restricciones:} Límites y reglas de qué NO hacer (Ej. No usar lenguaje técnico).
5. \textbf{Formato:} Resultado esperado (Tabla, correo de 2 párrafos, viñetas).

\end{frame}



\begin{frame}[fragile]{Ejercicio Práctico: Aplicando los 5 Pilares}


\textbf{Contexto:} Necesitas planear las vacaciones familiares, pero todos tienen gustos diferentes.
\textbf{Tarea:} Usar los 5 pilares para obtener un itinerario perfecto.

\begin{itemize}
  \item \textbf{Ejemplo de Prompt:}
\end{itemize}
\textit{``[ROL] Actúa como un agente de viajes experto. [CONTEXTO] Voy a viajar a Madrid por 4 días con mi familia (dos adultos y dos niños pequeños). Nos gusta la historia, pero los niños se aburren en los museos. [INSTRUCCIÓN] Crea un itinerario equilibrado. [RESTRICCIONES] No incluyas actividades que duren más de 2 horas seguidas. [FORMATO] Presenta el itinerario en una tabla organizada por días y horarios.''}

\end{frame}



\begin{frame}[fragile]{3. El Marco TAM (Tarea, Audiencia, Método)}


Para escribir correos o mensajes importantes, la estructura del \textit{Por qué} es tan importante como el \textit{Qué}.

\begin{itemize}
  \item \textbf{T - Tarea:} ¿Qué se necesita hacer exactamente? (Ej. Cancelar una suscripción, pedir un descuento).
  \item \textbf{A - Audiencia:} ¿Quién va a leer esto? ¿Cuál es su estado emocional o nivel de conocimiento? (Ej. Un representante de servicio al cliente saturado).
  \item \textbf{M - Método:} ¿Cómo se debe entregar la información? (Ej. Un correo amable, directo y estructurado con viñetas).
\end{itemize}
\end{frame}



\begin{frame}[fragile]{Ejercicio Práctico: El Marco TAM en Acción}


\textbf{Contexto:} Un albañil o contratista hizo mal una remodelación en tu casa y necesitas exigirle que la repare sin que se ofenda ni abandone la obra.

\begin{itemize}
  \item \textbf{Ejemplo de Prompt:}
\end{itemize}
Necesito enviar un mensaje por WhatsApp sobre un problema en la obra.
\textbf{Tarea:} Explicar que la pintura en la sala tiene manchas y debe repasarse este viernes sin costo extra.
\textbf{Audiencia:} El contratista a cargo de la obra (muy orgulloso de su trabajo, pero algo informal).
\textbf{Método:} Un mensaje amable pero firme, reconociendo primero el buen trabajo que hizo en la cocina, para luego pedirle que solucione la sala.

\end{frame}



\begin{frame}[fragile]{Recapitulando el Módulo}


\textbf{Hoy hemos aprendido:}
\begin{itemize}
  \item Dejar de usar palabras clave y empezar a dar contexto.
  \item Los \textbf{5 Pilares} garantizan resultados útiles y predecibles.
  \item El \textbf{Marco TAM} nos ayuda a adaptar nuestro mensaje a la realidad de quien lo recibe.
\end{itemize}
\textbf{¡Sigan aprendiendo, sigan trabajando y sigan disfrutando de la Inteligencia Artificial!}
\end{frame}




\begin{frame}[allowframebreaks]{References}
  \bibliographystyle{apalike}
  \bibliography{references}
\end{frame}


\end{document}